\subsection{Рівновіддалені вузли та феномен Рунге}

Інтуїтивний підхід до задачі наближення функцій часто підштовхує до вибору рівновіддалених точок дискретизації. Однак використання рівномірної сітки для побудови поліномів високого ступеня є вкрай нестабільним і навіть «катастрофічно поганим» методом. Математична проблема, що виникає при цьому, відома як феномен Рунге.

Це явище було відкрите німецьким математиком Карлом Рунге у 1901 році, коли він досліджував інтерполяцію ідеально гладкої аналітичної функції (яка тепер носить його ім'я): $f(x) = \frac{1}{1 + 25x^2}$. Рунге виявив парадоксальний результат: зі збільшенням кількості рівновіддалених вузлів інтерполяції $n$, максимальна похибка не зменшується, а навпаки — експоненціально зростає. Ближче до країв інтервалу інтерполяційний поліном починає здійснювати гігантські неконтрольовані осциляції, повністю відриваючись від значень цільової функції~\cite{trefethen2013approximation}. 

З точки зору математичного аналізу, феномен Рунге пояснюється поведінкою константи Лебега для рівномірних сіток. Доведено, що для рівновіддалених вузлів ця константа зростає з асимптотичною швидкістю $O(2^n / (n \log n))$. Це означає, що процес інтерполяції стає експоненціально чутливим до найменших змін або похибок округлення у вхідних даних.

У контексті цифрової обробки зображень феномен Рунге є критичною перешкодою. Цифрові сенсори та дисплеї фіксують і відтворюють зображення на ідеально рівномірній ортогональній сітці пікселів. Якби для зміни масштабу або просторового перетворення цілого зображення використовувалася глобальна поліноміальна інтерполяція за цими рівновіддаленими пікселями, це призвело б до повного руйнування візуальної інформації на краях зображення через екстремальні викиди яскравості та кольору.

Саме неможливість використання рівномірної сітки для стабільної апроксимації високого порядку фундаментально обґрунтовує необхідність переходу до нерівномірних сіток, або поліномів низького степеня. Як доводиться в теорії наближень, щоб нівелювати феномен Рунге та забезпечити рівномірну збіжність інтерполяційного процесу, вузли повинні бути асимптотично згущені на краях інтервалу. Ідеальною математичною моделлю, що задовольняє цю умову і гарантує алгоритмічну стійкість, є розподіл вузлів Чебишева. Таким чином, відмова від рівновіддалених вузлів на користь чебишовської сітки є не просто оптимізацією, а єдино коректним математичним шляхом для побудови стійких поліноміальних інтерполянтів.